
%\input{preamble1.tex}

% Document class
\documentclass[12pt]{article}

% Change font
\usepackage[T1]{fontenc}
\usepackage{lmodern}

% Text width
%\setlength{\columnsep}{15mm}
%\setlength{\columnseprule}{0.001mm}
%\addtolength{\oddsidemargin}{-0.5cm}
%\addtolength{\evensidemargin}{-0.5cm}
%\addtolength{\textwidth}{0.5cm}

%\addtolength{\topmargin}{-3cm}
%\addtolength{\textheight}{6cm}

\usepackage[margin=2cm]{geometry}
\usepackage{parskip}

%%% Packages %%%
	
	% Formatting
	\usepackage{setspace}
	\setstretch{1.1}
	
	%\usepackage{fontspec}
	%\setmainfont{Times New Roman}
	
	% Mathematics
	\usepackage{amsmath, latexsym}
	\usepackage{amsthm}
	\usepackage{amssymb}
	\usepackage{bm}
	\usepackage{eurosym}
	\usepackage{caption}
	\usepackage{accents}
	\usepackage{enumerate}
	\usepackage{empheq}
	\usepackage{array}
	\usepackage{booktabs}
	\usepackage{footmisc}
	\usepackage{dcolumn}
	\usepackage{verbatim}
	\usepackage{xcolor}
	\usepackage{tikz}
	\usepackage{bigints}
	\usetikzlibrary{patterns}
	%\usepackage{subcaption}
%	\usepackage[subrefformat=parens,labelformat=parens]{subcaption}
\usepackage[labelformat=simple]{subcaption}
\renewcommand\thesubfigure{(\alph{subfigure})}

    %\usepackage[capposition=top]{floatrow}
	% Define figure note command
\newcommand{\fignote}[1]{\caption*{\scriptsize Note: #1}}

	\usepackage{adjustbox}
	%\definecolor{mygray}{gray}{0.5}
	\usepackage{setspace} 
	\usepackage[flushleft]{threeparttable}	
	\usepackage[strict]{changepage}
	
	% Fonts
	\usepackage{dsfont}
	%\usepackage{pxfonts}
	\usepackage{mathrsfs}
	\usepackage{bbm}
	
	% Graphics
	\usepackage{graphicx}
	
	% For fixing location of figures
	\usepackage{float}
	
	% For landscape
	\usepackage{pdflscape}
	
	% Colors
	\definecolor{OliveGreen}{RGB}{53, 98, 26}
	\definecolor{LightOliveGreen}{RGB}{121,148,103}
	\definecolor{MyOrange}{RGB}{231,101,45}
	\definecolor{MyBlue}{RGB}{60,106,209}
	\definecolor{MyDarkOrange}{RGB}{169,74,1}

	% Bold and colored fonts
	\newcommand\boldmo[1]{\color{MyOrange}{\bm{#1}}}
	\newcommand\boldmb[1]{\color{MyBlue}{\bm{#1}}}
	\newcommand\boldmg[1]{\color{OliveGreen}{\bm{#1}}}
	
	% bold and colored text
	\newcommand\boldtg[1]{\textcolor{OliveGreen}{\textbf{#1}}}
	\newcommand\boldto[1]{\textcolor{MyOrange}{\textbf{#1}}}
	\newcommand\boldtb[1]{\textcolor{MyBlue}{\textbf{#1}}}
	
	% Hyperlinks
	\usepackage{hyperref}
%	\hypersetup{
%  colorlinks   = true ,    % Colours links instead of ugly 	boxes
%  urlcolor     = MyDarkOrange,    % Colour for external hyperlinks
%  linkcolor    = MyDarkOrange,%black,    % Colour of internal links
%  citecolor    = MyDarkOrange      % Colour of citations
%}
\hypersetup{
	colorlinks = true,
    linkcolor = MyDarkOrange, %blue
    urlcolor  = MyDarkOrange,
    citecolor = MyDarkOrange,
    anchorcolor = MyDarkOrange,
    hypertexnames=true
}

	% Appendix table of contents (loaded after hyperref)
	\usepackage{etoc}
	
	% Bibliography

%\usepackage{natbib}
\usepackage[backend=biber,style=authoryear,citestyle=apa,doi=false, eprint=false, url=false]{biblatex}
%\bibliographystyle{aeanobold-oxford}

%%% Definitions of shortcuts %%%

	% Equations
	\def\dsp{\displaystyle}
	\def\beq{\begin{equation}}
	\def\eeq{\end{equation}}
	\def\ben{\begin{equation*}}
	\def\een{\end{equation*}}
	\def\ba{\begin{eqnarray}}
	\def\ea{\end{eqnarray}}
	\def\ban{ \begin{eqnarray*} }
	\def\ean{ \end{eqnarray*} }
	\def\bea{\begin{array}}
	\def\eea{\end{array}}
	
	% Itemize and enumerate
	\def\bi{\begin{itemize}}
	\def\ei{\end{itemize}}

	\def\bnum{\begin{enumerate}}
	\def\enum{\end{enumerate}}
	\def\bnumi{\begin{enumerate}[(i)]}


	% Greek letters
	\def\k{\kappa}
	\def\r{\rho}
	\def\a{\alpha}
	\def\b{\beta}
	\def\g{\gamma}
	\def\gh{\hat\g}
	\def\G{\Gamma}
	\def\d{\delta}
	\def\D{\Delta}
	\def\e{\varepsilon}
	\def\f{\phi}
	\def\F{\Phi}
	\def\p{\psi}
	\def\pb{\bar\p}
	\def\P{\Psi}
	\def\C{\Chi}
	\def\Th{\Theta}
	\def\l{\lambda}
	\def\m{\mu}
	\def\n{\nu}
	\def\o{\omega}
	\def\O{\Omega}
	\def\l{\lambda}
	\def\L{\Lambda}
	\def\s{\sigma}
	\def\S{\Sigma}
	\def\vf{\varphi}
	\def\vphi{\varphi}
	\def\t{\theta}
	\def\ph{\varphi}
	\def\vs{\varsigma}
	\def\z{\zeta}

	% Caligraphic letters
	\def\cN{{\cal N}}
	\def\cM{{\cal M}}
	\def\cL{{\cal L}}
	\def\cF{{\cal F}}
	\def\cA{{\cal A}}
	\def\cB{{\cal B}}
	\def\cC{{\cal C}}
	\def\cO{{\cal O}}
	\def\cD{{\cal D}}
	\def\cR{{\cal R}}
	\def\cQ{{\cal Q}}
	\def\cG{{\cal G}}
	\def\cP{{\cal P}}
	\def\cT{{\cal T}}
	\def\cH{{\cal H}}
	\def\cU{{\cal U}}
	\def\cJ{{\cal J}}
	\def\cI{{\cal I}}
	\def\cH{{\cal H}}
	\def\cS{{\cal S}}
	\def\cZ{{\cal Z}}
	\def\cX{{\cal X}}
	\def\cE{{\cal E}}
	\def\cV{{\cal V}}
	\def\cK{{\cal K}}
	\def\cW{{\cal W}}
	
	% Operators
	\DeclareMathOperator*{\argmin}{argmin}
	\DeclareMathOperator*{\argmax}{argmax}
	
	% Properties, Lemmas et theorems...
	\theoremstyle{plain} \newtheorem{prop}{Proposition}
	\theoremstyle{plain} \newtheorem{assum}{Assumption}
	\theoremstyle{plain} \newtheorem{thm}{Theorem}
	\theoremstyle{plain} \newtheorem{lemma}{Lemma}
	\theoremstyle{plain} \newtheorem{cor}{Corollary}
	\theoremstyle{plain} \newtheorem{defi}{Definition}
   	\theoremstyle{plain} \newtheorem{claim}{Claim}
	\theoremstyle{plain} \newtheorem{rem}{Remark}
	\theoremstyle{definition} \newtheorem{def-prop}{Def-Prop}

	\def\bdefi{\begin{defi}}
	\def\edefi{\end{defi}}
	\def\bprop{\begin{prop}}
	\def\eprop{\end{prop}}
	\def\blem{\begin{lemma}}
	\def\elem{\end{lemma}}
	\def\bassu{\begin{assum}}
	\def\eassu{\end{assum}}
	\def\bthm{\begin{thm}}
	\def\ethm{\end{thm}}
	\def\bth{\begin{thm}}
	\def\eth{\end{thm}}
	\def\bcor{\begin{cor}}
	\def\ecor{\end{cor}}
	\def\brem{\begin{rem}}
	\def\erem{\end{rem}}
	\def\bdefprop{\begin{def-prop}}
	\def\edefprop{\end{def-prop}}
	\def\bclaim{\begin{claim}}
	\def\eclaim{\end{claim}}



	% Thick math sets
	\def\1{\mathds{1}}
	\def\P{\mathbb{P}}
	\def\R{\mathbb{R}}
	\def\C{\mathbb{C}}
	\def\E{\mathbb{E}}
	\def\N{\mathbb{N}}
	\def\Q{\mathbb{Q}}
	\def\Sb{\mathbb{S}}

	% Convergence
	\def\ascv{\overset{a.s.}{\longrightarrow}}
	\def\dcv{\Longrightarrow}
	\def\pcv{\overset{P}{\longrightarrow}}

	% Economics and econometrics
	\def\rel{\succeq}
	\def\srel{\succ}
	\def\hb{\hat{\b}}
	\def\hbo{\hat{\b}_{OLS}}
	\def\hth{\hat{\t}}
	\def\ut{\underline{\t}}

	\def\bx{\overline{x}}
	\def\bo{\overline{\o}}
	\def\by{\overline{y}}
	
	% Partial derivative
	\newcommand\parder[2]{\frac{\partial #1}{\partial #2}}

%%% Others %%%

	% Allow large matrices
	\setcounter{MaxMatrixCols}{20}
	
%%% Specific %%%
\def\hb{\overline{h}}

\def\vp{\varphi}
\def\uvp{\underline{\varphi}}
\def\uw{\underline{w}}
\def\os{\underline{s}}
\def\ovp{\overline{\vp}}
\def\hvp{\hat{\vp}}
\def\tvp{\tilde{\vp}}
\def\op{\overline{\phi}}
\def\up{\underline{\phi}}

\def\al{\underline{a}}
\def\au{\overline{a}}


\def\rl{\underline{r}}
\def\rh{\overline{r}}

\def\zl{\underline{z}}
\def\zu{\overline{z}}
\def\xl{\underline{x}}
\def\xu{\overline{x}}
\def\yl{\underline{y}}
\def\yu{\overline{y}}
\def\sl{\underline{s}}
\def\su{\overline{s}}
\def\nul{\underline{\nu}}
\def\nuu{\overline{\nu}}
\def\bcF{\overline{\cF}}
\def\hl{\underline{h}}
\def\hu{\overline{h}}

\def\cd{\cdot}



\newcommand{\x}{\bm{x}}
\newcommand{\y}{\bm{y}}
\newcommand{\h}{\overline{h}}

\newcommand{\vl}{\underline{v}}
\newcommand{\vu}{\overline{v}}

\newcommand{\wl}{\underline{w}}
\newcommand{\wu}{\overline{w}}


\newcommand{\zetal}{\underline{\zeta}}
\newcommand{\zetau}{\overline{\zeta}}

\newcommand{\elll}{\underline{\ell}}
\newcommand{\ellu}{\overline{\ell}}

\newcommand{\gl}{\underline{g}}
\newcommand{\gu}{\overline{g}}

\newcommand{\bl}{\underline{b}}
\newcommand{\bu}{\overline{b}}

%\renewcommand*{\thefootnote}{\fnsymbol{footnote}}

\newcommand{\smallinput}[1]{
  \begingroup
  \renewcommand\normalsize{\small}
  \input{#1}
  \endgroup
}

\newcommand{\footnotesizeinput}[1]{
  \begingroup
  \renewcommand\normalsize{\footnotesize}
  \input{#1}
  \endgroup
}

\newcommand{\scriptsizeinput}[1]{
  \begingroup
  \renewcommand\normalsize{\scriptsize}
  \input{#1}
  \endgroup
}

\newcommand{\tinyinput}[1]{
  \begingroup
  \renewcommand\normalsize{\tiny}
  \input{#1}
  \endgroup
}


% Colored text
\newcommand{\gray}{\textcolor{gray}}
\newcommand{\tgray}[1]{\text{\textcolor{gray}{#1}}}


% Stuff I am not happy about
\newcommand{\todo}[1]{\textbf{\textcolor{red}{#1}}}


% Inner product
\newcommand{\la}{\left\langle}
\newcommand{\ra}{\right\rangle}

% Acrsinh functions etc
\DeclareMathOperator{\arccosh}{arcCosh}
\DeclareMathOperator{\arcsinh}{arcsinh}

% bold
\newcommand{\bg}{\pmb{g}}
\newcommand{\bh}{\pmb{h}}

\addbibresource{biblio.bib}

\title{\textbf{Geographic Markups: Notes and Derivations}}
\author{Santiago Franco}
\date{\today}

% Scratch paper for exploratory derivations.
% Conventions:
%   - One \section per topic. Start each new entry with a dated comment.
%   - Flag non-rigorous steps and open questions inline with \textbf{Open:}.
%   - When a result is settled and moved to Paper.tex, leave a one-line note saying where it went.

\begin{document}
\maketitle


%%%%%%%%%%%%%%%%%%%%%%%%%%%%%%%%%%%%%%%%%%%%%%%%%%%%%%%%%%%%%%%%%%%%%%%%%%%%%%%%%%%%%%%%%
%									baseline framework									%
%%%%%%%%%%%%%%%%%%%%%%%%%%%%%%%%%%%%%%%%%%%%%%%%%%%%%%%%%%%%%%%%%%%%%%%%%%%%%%%%%%%%%%%%%

% 2026-11-24: first draft, adapted from the multi-sector appendix of the JMP (a:muext).
% 2026-11-24: revised. Regions n, industries i, time t; no outside goods; measurement error in observed sales;
%             fixed inputs and arbitrary returns to scale; input prices specific to region and industry.

\section{Baseline Framework}\label{sec:baseline}

\subsection{Set Up}\label{sec:setup}

Time is discrete and indexed by $t$.
There are $N$ regions indexed by $n \in \{1, \dots, N\}$.
Regions differ in their population, $L_{nt}$.
There are $I$ industries indexed by $i \in \{1, \dots, I\}$.
Each industry produces non-traded varieties, so varieties produced in region $n$ are consumed only in region $n$.
Let $\cJ_{int}$ denote the finite set of establishments active in industry $i$, region $n$, and period $t$, indexed by $j$.
Each establishment produces a single differentiated variety.
The set $\cJ_{int}$ is taken as given.

All decisions below are static, so the model is a sequence of period-$t$ equilibria.
Time enters through changes in populations, incomes, the sets of establishments, productivity, and input prices.
The only intertemporal element is that some inputs are chosen before period $t$ (see Section~\ref{sec:production}).

\subsection{Preferences}\label{sec:preferences}

Each region is populated by $L_{nt}$ identical households.
Each household has income $E_{nt}$, which it takes as given and spends entirely on local goods.
Households have Cobb-Douglas preferences over industry bundles,
\begin{equation}\label{eq:uppertier}
	U_{nt} = \prod_{i=1}^{I} \left(\frac{Y_{int}}{\theta_i}\right)^{\theta_i}, \qquad \sum_{i=1}^{I} \theta_i = 1,
\end{equation}
where $Y_{int}$ is per-capita consumption of the industry-$i$ bundle in region $n$ and period $t$.

Within each industry, households have HDIA preferences over varieties, in the terminology of \textcite{Matsuyama&Ushev17}.
These are the Kimball preferences of \textcite{Kimball95}.
The industry bundle $Y_{int}$ is implicitly defined by
\begin{equation}\label{eq:hdia}
	\sum_{j \in \cJ_{int}} \Upsilon_i\left(\frac{y_{jint}}{Y_{int}}\right) = 1,
\end{equation}
where $y_{jint}$ is per-capita consumption of variety $j$.

\bassu[HDIA aggregator]\label{assum:hdia}
For each industry $i$, the aggregator $\Upsilon_i(\cdot)$ is twice continuously differentiable, strictly increasing, and strictly concave, with $\Upsilon_i(0) = 0$.
\eassu

CES is the special case $\Upsilon_i(q) = q^{(\sigma_i - 1)/\sigma_i}$.
The aggregator is industry specific, so the curvature of demand, and hence markups, can differ across industries.
The preference parameters $\theta_i$ and $\Upsilon_i$ are constant over time and across regions.

\textbf{Upper tier.}
Because both tiers are homothetic, each industry bundle has a price index $\P_{int}$, and the household maximizes \eqref{eq:uppertier} subject to $\sum_i \P_{int} Y_{int} = E_{nt}$.
Cobb-Douglas implies constant expenditure shares:
\begin{equation}\label{eq:indexp}
	\P_{int} Y_{int} = \theta_i E_{nt}.
\end{equation}

\textbf{Lower tier.}
For a given $Y_{int}$, the household chooses $\{y_{jint}\}_{j}$ to minimize expenditure $\sum_{j} p_{jint} y_{jint}$ subject to \eqref{eq:hdia}.
With Lagrange multiplier $\lambda_{int}$, the first-order condition for variety $j$ is
\begin{equation*}
	p_{jint} = \frac{\lambda_{int}}{Y_{int}} \Upsilon_i'\left(\frac{y_{jint}}{Y_{int}}\right).
\end{equation*}
Define the competition price index $\mathbb{D}_{int} \equiv \lambda_{int}/Y_{int}$ and the inverse marginal aggregator $\vp_i(\cdot) \equiv (\Upsilon_i')^{-1}(\cdot)$.
Strict concavity of $\Upsilon_i$ makes $\Upsilon_i'$ strictly decreasing, so $\vp_i$ is well defined and strictly decreasing.
The per-capita demand for variety $j$ is then
\begin{equation}\label{eq:demand}
	\frac{y_{jint}}{Y_{int}} = \vp_i\left(\frac{p_{jint}}{\mathbb{D}_{int}}\right).
\end{equation}
Substituting \eqref{eq:demand} into \eqref{eq:hdia}, the competition price index is implicitly defined by
\begin{equation}\label{eq:Dindex}
	\sum_{j \in \cJ_{int}} \Upsilon_i\left(\vp_i\left(\frac{p_{jint}}{\mathbb{D}_{int}}\right)\right) = 1.
\end{equation}
Equation \eqref{eq:Dindex} pins down $\mathbb{D}_{int}$ from prices alone, so $\mathbb{D}_{int}$ does not depend on $Y_{int}$.
Since minimized expenditure is $\sum_j p_{jint} y_{jint} = \P_{int} Y_{int}$, dividing by $Y_{int}$ and using \eqref{eq:demand} gives the ideal price index
\begin{equation}\label{eq:Pindex}
	\P_{int} = \sum_{j \in \cJ_{int}} p_{jint} \, \vp_i\left(\frac{p_{jint}}{\mathbb{D}_{int}}\right).
\end{equation}
Under CES, $\P_{int}$ and $\mathbb{D}_{int}$ are proportional.
Outside CES they are not, and the ratio $\P_{int}/\mathbb{D}_{int}$ will matter for markup estimation.

\textbf{Market demand.}
Total demand for variety $j$ is $q_{jint} = L_{nt} y_{jint}$.
Combining \eqref{eq:indexp} and \eqref{eq:demand},
\begin{equation}\label{eq:marketdemand}
	q_{jint} = \frac{\theta_i E_{nt} L_{nt}}{\P_{int}} \, \vp_i\left(\frac{p_{jint}}{\mathbb{D}_{int}}\right).
\end{equation}
The price elasticity of demand, holding the price indices fixed, depends only on the relative price:
\begin{equation}\label{eq:sigma}
	\s_i\left(\frac{p_{jint}}{\mathbb{D}_{int}}\right) \equiv -\frac{\partial \log q_{jint}}{\partial \log p_{jint}} = -\frac{\frac{p_{jint}}{\mathbb{D}_{int}} \vp_i'\left(\frac{p_{jint}}{\mathbb{D}_{int}}\right)}{\vp_i\left(\frac{p_{jint}}{\mathbb{D}_{int}}\right)}.
\end{equation}

\subsection{Production}\label{sec:production}

Establishment $j$ in industry $i$, region $n$, and period $t$ produces with a Hicks-neutral Cobb-Douglas technology over a set of inputs $\cX$:
\begin{equation}\label{eq:prodfn}
	q_{jint} = \o_{jint} \prod_{x \in \cX} \left(v^x_{jint}\right)^{\b^x_i},
\end{equation}
where $\o_{jint}$ is establishment productivity, $v^x_{jint}$ is the quantity of input $x$, and $\b^x_i \geq 0$ is the output elasticity of input $x$.
Output elasticities are common to all establishments in industry $i$ and constant over time.
An input not used in industry $i$ has $\b^x_i = 0$.
Returns to scale, $\g_i \equiv \sum_{x \in \cX} \b^x_i$, are not restricted.

\textbf{Inputs.}
Input markets are competitive.
Establishments in industry $i$ and region $n$ take the price of input $x$ in period $t$, $W^x_{int}$, as given.
Inputs are of two types.
A subset $\cX^f \subseteq \cX$ is flexible: these inputs are chosen in period $t$ after $\o_{jint}$ is known, and can be freely adjusted.
The remaining inputs, $\cX \setminus \cX^f$, are fixed: they are chosen before period $t$ and are given within the period.
Labor, $x = l$, is flexible, with $v^l_{jint} = l_{jint}$ and $W^l_{int} = W_{int}$.

\textbf{Market structure.}
Establishments compete under monopolistic competition.
Each establishment takes the price indices $\P_{int}$ and $\mathbb{D}_{int}$, industry expenditure $\theta_i E_{nt} L_{nt}$, and input prices as given.
With a finite set $\cJ_{int}$, an establishment's price does affect $\mathbb{D}_{int}$ and $\P_{int}$ through \eqref{eq:Dindex} and \eqref{eq:Pindex}.
Monopolistic competition assumes establishments do not internalize this effect.
The oligopoly extension in the style of \textcite{Atkeson&Burstein08} will relax this.

\subsection{Markups}\label{sec:markups}

\textbf{Cost minimization.}
Fix output $q_{jint}$ and the levels of the fixed inputs.
For a flexible input $x \in \cX^f$, the first-order condition of cost minimization is
\begin{equation*}
	W^x_{int} = \l_{jint} \, \b^x_i \frac{q_{jint}}{v^x_{jint}},
\end{equation*}
where $\l_{jint}$, the multiplier on the output constraint, is short-run marginal cost.
Define the markup as price over marginal cost, $\m_{jint} \equiv p_{jint}/\l_{jint}$, and the expenditure share of input $x$ in true sales as $\a^x_{jint} \equiv W^x_{int} v^x_{jint} / (p_{jint} q_{jint})$.
Multiplying the first-order condition by $v^x_{jint}/(p_{jint} q_{jint})$ and rearranging gives the ratio estimator of \textcite{DeLoecker&Warzynski12}:
\begin{equation}\label{eq:ratio}
	\m_{jint} = \frac{\b^x_i}{\a^x_{jint}}, \qquad x \in \cX^f.
\end{equation}
This step uses only cost minimization over a flexible, price-taking input.
It does not depend on the demand system, the market structure, returns to scale, or the fixed inputs.

\textbf{Pricing.}
Given the short-run cost function, the establishment chooses $p_{jint}$ to maximize profits subject to \eqref{eq:marketdemand}, taking the price indices as given.
The first-order condition is the Lerner condition:
\begin{equation}\label{eq:lerner}
	\m_{jint} = \m_i\left(\frac{p_{jint}}{\mathbb{D}_{int}}\right) \equiv \frac{\s_i\left(p_{jint}/\mathbb{D}_{int}\right)}{\s_i\left(p_{jint}/\mathbb{D}_{int}\right) - 1}.
\end{equation}
This requires $\s_i > 1$ at the optimum, so establishments operate on the elastic part of their demand curve.
The Lerner condition holds whatever the shape of marginal cost, so it does not depend on returns to scale or the fixed inputs either.

\textbf{Open:} Marginal cost varies with output unless short-run returns to scale, $\g^f_i \equiv \sum_{x \in \cX^f} \b^x_i$, equal one.
The second-order condition for pricing then needs joint conditions on $\Upsilon_i$ and $\g^f_i$.
These are not stated yet.

\textbf{Measurement error.}
Sales are observed with error.
Observed sales are
\begin{equation}\label{eq:obssales}
	\tilde{R}_{jint} = p_{jint} q_{jint} \, e^{\n_{jint}},
\end{equation}
where $\n_{jint}$ is measurement error.
Input expenditures $W^x_{int} v^x_{jint}$ are observed without error.
The observed expenditure share is $\tilde{\a}^x_{jint} \equiv W^x_{int} v^x_{jint}/\tilde{R}_{jint} = \a^x_{jint} e^{-\n_{jint}}$.
Substituting into \eqref{eq:ratio},
\begin{equation}\label{eq:ratioobs}
	\m_{jint} = \frac{\b^x_i}{\tilde{\a}^x_{jint}} \, e^{-\n_{jint}}, \qquad x \in \cX^f.
\end{equation}
Because the error is in measurement only, establishments' choices are unaffected and \eqref{eq:lerner} still holds at the true price.

\textbf{Open:} The sales shares used later are built from observed sales, so they contain $\n_{jint}$ as well.
This creates a correlation between the regressors and the error in the estimating equation.
It motivates the lagged-sales-share instruments used in the JMP.
The assumptions needed on $\n_{jint}$ (mean zero, independence across establishments and over time) are not stated yet.

\textbf{Summary.}
Equations \eqref{eq:ratioobs} and \eqref{eq:lerner} are the two building blocks for markup estimation.
The first links markups to observables up to $\b^x_i$ and the measurement error.
The second links markups to the demand system through the relative price $p_{jint}/\mathbb{D}_{int}$, which is not observed.
The next step expresses the relative price in terms of observed sales shares, as in Appendix \texttt{a:muestderiv} of the JMP.


\end{document}
